\documentclass{article}
\usepackage{amsmath, amssymb, amsthm}
\newtheorem{proposition}{Proposition}
\newtheorem{theorem}{Theorem}
\newtheorem{corollary}{Corollary}
\newtheorem{remark}{Remark}
\usepackage{graphicx}
\usepackage{hyperref}
\usepackage{geometry}
\usepackage{algorithm}
\usepackage{algpseudocode}
\usepackage{xcolor}
\usepackage{float}
\geometry{margin=1in}

\newcommand{\merw}{EigenTree}
\title{\merw: A Decentralized Communication Protocol\\Achieving Centralized Multi-Armed Bandit Results}
\author{Salvatore Rastelli, Andrea Perin}
\date{May 2026}

\begin{document}

\maketitle

\begin{abstract}
	In decentralized cooperative learning, no agent has global knowledge, and yet centralized coordination is statistically efficient.
	We propose \merw{}, a decentralized communication protocol on undirected graphs that establishes, from purely local interactions, a de-facto coordinator among the agents.
	\merw{} uses eigenvector centrality---as computed by a gossip-based power iteration---to 1) choose the most central node as coordinator, and 2) construct a routing tree, rooted at the coordinator node, which spans the graph.
	Information flows according to this routing tree: from leaves to root during the aggregation phase, and vice versa during the boradcasting phase.
	Thus, a centralized aggregator emerges out of purely decentralized interactions, at the cost of introducing appropriate communication rounds between agents in the tree.
	The resulting protocol can be directly used in various decentralized cooperative learning setups, of which we show two---regret minimization and best arm identification---where the method recovers centralized-like results.
\end{abstract}

\section{Introduction}

In the cooperative multi-agent multi-armed bandit (MA-MAB) setting, $N$ agents are connected by an undirected graph $G = (V, E)$. At each round, every agent pulls one arm and receives a stochastic reward; agents communicate with neighbors to improve collective performance. Two objectives are studied: \emph{regret minimization}, where the goal is to minimize cumulative group regret over a horizon $T$; and \emph{Best Arm Identification} (BAI), where the goal is to identify the best arm with high confidence using as few total arm pulls as possible.

The central difficulty in both settings is the tension between \emph{decentralization} (no agent has global knowledge) and \emph{statistical efficiency} (centralized coordination is optimal). Existing approaches resolve this either by symmetric averaging (e.g., Coop-UCB2 \cite{landgren2021distributed}, which uses doubly stochastic consensus weights to diffuse information uniformly) or by assuming all-to-all communication (e.g., Hillel et al.\ \cite{hillel2013distributed}, whose distributed successive elimination requires every agent to broadcast to all others in each round).

\merw{} takes a different route: we show that a coordinator can \emph{emerge} from the graph structure itself, without any agent knowing the global topology. The key insight is that the Maximal Entropy Random Walk (MERW) eigenvector centrality, computed via a fully decentralized gossip protocol, identifies the structurally most central node. This node, the hub, then organizes the rest of the network into a spanning routing tree via max-flooding, again using only local information.

Once the tree is in place, the hub acts as a de-facto central coordinator. Every agent routes its observations up toward the hub, which sees the globally aggregated picture and makes decisions on behalf of the network. The downlink carries only what is needed: in the regret setting, the hub's enriched estimates; in the BAI setting, the surviving arm set and a timing signal so all agents restart each elimination round simultaneously.

The contribution of \merw{} is therefore twofold. First, it provides a fully decentralized protocol for \emph{inducing centrality}: the hub and tree emerge from local MERW computations, requiring no privileged node, no global communication, and no prior knowledge of $N$ or the graph structure. Second, it demonstrates that this emergent coordinator is statistically sufficient: the hub achieves centralized UCB performance for regret minimization, and the same sample complexity as Hillel's protocol for BAI, while the communication overhead over Hillel is only $O(D)$ extra time steps per elimination round, where $D$ is the depth of the MERW-induced tree.

\section{Related Work}
\label{sec:related}

\paragraph{Cooperative multi-agent bandits.}
The cooperative MA-MAB problem has been studied under two broad communication models. In the \emph{gossip/consensus} model, agents exchange information with neighbors and update local estimates via averaging; no coordinator is assumed. Coop-UCB2 \cite{landgren2021distributed} uses doubly stochastic consensus weights to diffuse observations across the graph, achieving sublinear group regret that depends on the spectral gap of the communication matrix. Martinez-Rubio et al.\ \cite{martinez2019decentralized} study a similar setting and introduce the total-pull clock as a way to index UCB confidence intervals by actual observations rather than wall-clock rounds, an idea we adopt directly. These consensus-based approaches are symmetric: every agent plays an equivalent role, and there is no notion of a coordinator.

Hillel et al.\ \cite{hillel2013distributed} study cooperative BAI under a fully connected communication model, where in each round every agent may broadcast a message to all others. Their main result is a tradeoff between communication rounds and sample complexity: with a single communication round, the best achievable speedup is $\sqrt{k}$ (and this is tight by their lower bound); with $O(\log(1/\varepsilon))$ rounds, their Algorithm 3 achieves the ideal factor-$k$ speedup. We build on Algorithm 3: \merw-BAI uses the same pull schedule and elimination rule, but replaces the all-to-all broadcast with hop-by-hop relay over the MERW spanning tree. Because every agent receives the same global information after each broadcast in Hillel's setting, all agents compute the same surviving arm set independently; in our setting, the hub is the only node that sees the full aggregate and computes the elimination -- a role that emerges from the graph structure rather than being assumed by the communication model.

\paragraph{Eigenvector centrality and MERW.}
Eigenvector centrality scores nodes by their contribution to the leading eigenvector of the adjacency matrix and is widely used in network analysis. The Maximal Entropy Random Walk \cite{burda2009localization} provides an information-theoretic grounding for this measure: $\psi_i$ is the unique centrality consistent with assigning equal probability to all paths of equal length. The localization phenomenon -- that $\psi_i^2$ concentrates sharply on a small set of nodes in heterogeneous graphs -- was characterized by Burda et al.\ \cite{burda2009localization} and motivates using $\psi$ to identify a hub rather than, say, degree or betweenness centrality.

\paragraph{Decentralized eigenvector computation.}
Distributed power iteration for computing the leading eigenvector of a matrix defined by a network has been studied in the systems literature. Jelasity et al.\ \cite{jelasity2007asynchronous} give a gossip-based normalization that avoids the global all-reduce required by standard power iteration, at the cost of $O(\log N)$ extra gossip rounds per iteration. We use their protocol directly; the key observation for our setting is that gossip normalization never requires an agent to know $N$, which is consistent with our fully decentralized model.

\section{MERW and Hub Identification}
\label{sec:merw}

\subsection{Maximal Entropy Random Walk}

The Maximal Entropy Random Walk \cite{burda2009localization} is the unique stationary Markov chain on $G$ that assigns equal probability to all paths of equal length between any two fixed endpoints. It is the maximally unbiased random walk consistent with the graph structure.

Let $A \in \{0,1\}^{N \times N}$ be the adjacency matrix of $G$. By the Perron--Frobenius theorem, $A$ has a unique positive leading eigenvalue $\lambda_1$ with a corresponding strictly positive eigenvector $\psi \in \mathbb{R}^N_{>0}$, normalized so that $\|\psi\|_2 = 1$. The MERW transition matrix is
\[
S_{ij} = \frac{A_{ij}}{\lambda_1} \cdot \frac{\psi_j}{\psi_i},
\]
and its unique stationary distribution is $\pi_i = \psi_i^2$. Unlike the standard random walk, whose stationary distribution is proportional to degree, MERW concentrates its stationary mass on structurally central nodes. We use $\psi_i$ as the centrality score of node $i$; see Appendix~\ref{app:merw} for a detailed discussion of the localization phenomenon and its information-theoretic justification.

\subsection{Decentralized Hub Identification}

In a multi-agent setting, no single agent has access to the full adjacency matrix $A$. We compute $\psi$ via the distributed power iteration of \cite{jelasity2007asynchronous}, in which each agent normalizes using only local gossip rather than a global reduce. Crucially, the gossip normalization replaces the standard all-reduce with pairwise averaging of log-growth rates, so no agent ever needs to know $N$ or the global topology. Since we require only the \emph{ratio} $\psi_j / \psi_i$ to determine routing directions, and not the absolute eigenvector values, the unnormalized iterates suffice; see Appendix~\ref{app:poweriter} for the full protocol and convergence analysis.

After power iteration, each agent $i$ holds $\tilde{\psi}_i \approx c \cdot \psi_i$ for some common unknown scale $c > 0$. A short \emph{max-flooding} phase propagates the global maximum across the graph: each agent relays the largest $\tilde{\psi}$ value it has seen, and every local maximum re-routes toward the incoming direction. The result is a single directed spanning tree $\mathcal{T}$ rooted at the unique global maximum $h^\star = \arg\max_i \tilde{\psi}_i$, with all edges oriented from children toward the root. Convergence is guaranteed in at most $\mathrm{diam}(G)$ rounds; see Appendix~\ref{app:flooding} for the detailed protocol and stopping criterion.

\begin{proposition}[Centralized aggregation via routing]
\label{prop:aggregation}
Suppose each node $i \in V$ holds a scalar $x_i$ and propagates it toward hub $h^\star$ one hop per round, with each intermediate node forwarding the sum of what it receives. After $D$ rounds, the hub holds $\sum_{i \in V} x_i$.
\end{proposition}

\begin{proof}
By induction on tree depth. Leaf nodes (depth $D$) send $x_i$ to their parents after round 1. At each subsequent round, every node at depth $d$ receives the partial sums from all its children and forwards the total to its parent. After $D$ rounds the hub has accumulated all contributions.
\end{proof}

This aggregation property is the bridge between the decentralized routing structure and centralized statistical performance.

\begin{proposition}[Broadcast completeness]
\label{prop:broadcast}
Suppose the hub $h^\star$ holds a value $v$ and sends it one hop per round toward its children, with each intermediate node forwarding the received value to all its children. After $D$ rounds, every node $i \in V$ holds $v$.
\end{proposition}

\begin{proof}
By induction on depth. Nodes at depth 1 (direct children of $h^\star$) receive $v$ after round 1. If every node at depth $d-1$ holds $v$ after $d-1$ rounds, then in round $d$ each of them forwards $v$ to all their children (depth $d$). After $D$ rounds all nodes at depth $D$ have received $v$, completing the induction.
\end{proof}

\begin{theorem}[Coordinator sufficiency]
\label{thm:coordinator}
Let $\mathcal{A}$ be any bandit algorithm whose decision at each step depends only on the global sufficient statistics $\{(\hat{n}^k, \hat{s}^k)\}_{k=1}^K$ and the total pull count $P = \sum_{i \in V}\sum_{k=1}^K n_i^k$, where $\hat{n}^k = \sum_{i \in V} n_i^k$ and $\hat{s}^k = \sum_{i \in V} s_i^k$. Suppose a coordinator $h^\star$ exists that:
\begin{enumerate}
    \item receives exact global sufficient statistics from all agents within $D$ rounds (Proposition~\ref{prop:aggregation}), and
    \item computes $P$ by counting the uplink messages received each cycle plus its own pull, and broadcasts $(\hat{n}, \hat{s}, P)$ back to all agents within $D$ rounds (Proposition~\ref{prop:broadcast}), and
    \item all agents begin the next decision round simultaneously (synchronization via the depth signal $D$).
\end{enumerate}
Then the coordinator's information state at each decision point is identical to that of a centralized instance of $\mathcal{A}$. The coordinator therefore achieves the same statistical performance as $\mathcal{A}$ run on the pooled observations of all $N$ agents, with an additive communication overhead of at most $2D$ time steps per decision round.
\end{theorem}

\begin{proof}
Fix any decision round $r$. By Proposition~\ref{prop:aggregation}, after $D$ uplink steps the hub holds $\hat{n}^k = \sum_i n_i^k$ and $\hat{s}^k = \sum_i s_i^k$ for every arm $k$. The hub also counts the number of uplink messages received and adds one for its own pull to obtain $P = \sum_i \sum_k n_i^k$, the exact total number of arm pulls made by all agents in all rounds up to $r$. Together $(\hat{n}, \hat{s}, P)$ are exactly the sufficient statistics a centralized agent would maintain after observing the same pulls. Because $\mathcal{A}$'s decision depends only on these statistics, the hub computes the same action as the centralized agent. By Proposition~\ref{prop:broadcast}, after $D$ downlink steps every agent holds $(\hat{n}, \hat{s}, P)$ and the hub's decision. By the depth signal, all agents begin the next pull phase at the same time step. Therefore, in terms of arm pulls, the protocol is a faithful simulation of $\mathcal{A}$ on the pooled observations, and any performance guarantee for $\mathcal{A}$ carries over. The only cost is $2D$ idle time steps per round for communication.
\end{proof}

Figures~\ref{fig:merw_viz_ba}--\ref{fig:merw_viz_grid} illustrate the MERW centrality and induced routing tree on three representative graph families. On the Barab\'{a}si--Albert graph, $\psi^2$ concentrates sharply on the highest-degree hub. On the Erd\H{o}s--R\'{e}nyi graph, centrality is more spread but still identifies a clear hub. On the grid, MERW correctly selects the center node. Additional graph families (barbell, lollipop, chain, star, clustered) are shown in Appendix~\ref{app:merw_viz}.

\begin{figure}[H]
    \centering
    \includegraphics[width=\textwidth]{results/MERW_visualization/merw_tree_ba_N20.pdf}
    \caption{MERW centrality and routing tree on a Barab\'{a}si--Albert graph ($N=20$). Three panels: \emph{left} -- undirected graph with nodes colored by $\psi_i^2$; \emph{center} -- routing subtree overlaid on the graph, hub ringed in gold; \emph{right} -- induced directed tree in hierarchical layout, hub at top, arrows pointing child $\to$ parent.}
    \label{fig:merw_viz_ba}
\end{figure}

\begin{figure}[H]
    \centering
    \includegraphics[width=\textwidth]{results/MERW_visualization/merw_tree_er_N20.pdf}
    \caption{MERW centrality and routing tree on an Erd\H{o}s--R\'{e}nyi graph ($N=20$). Panel layout as in Figure~\ref{fig:merw_viz_ba}.}
    \label{fig:merw_viz_er}
\end{figure}

\begin{figure}[H]
    \centering
    \includegraphics[width=\textwidth]{results/MERW_visualization/merw_tree_grid_N25.pdf}
    \caption{MERW centrality and routing tree on a grid graph ($N=25$). Panel layout as in Figure~\ref{fig:merw_viz_ba}.}
    \label{fig:merw_viz_grid}
\end{figure}

\subsection{The Two Versions of \merw}

The hub and tree are constructed identically for both the regret and BAI objectives. We describe each version in the following sections.

\section{Regret Minimization}
\label{sec:regret}

\subsection{Protocol}

There are $K$ arms with unknown means $\mu_1 \geq \mu_2 \geq \cdots \geq \mu_K$ and sub-Gaussian reward distributions with parameter $\sigma^2$. The objective is to minimize the expected cumulative group regret:
\[
R(T) = \sum_{i \in V} \sum_{t=1}^{T} \mathbb{E}[\mu_1 - \mu_{a_i(t)}].
\]

Each relay cycle lasts exactly $2D + 1$ rounds. At the start of each cycle all $N$ agents pull simultaneously; the hub aggregates observations via the uplink, runs UCB on the global statistics, and broadcasts its estimates back down. Non-hub nodes update their local estimates upon receiving the downlink and schedule their next pull to align with the start of the next cycle. The total-pull count $P$ (indexed by actual arm pulls, not wall-clock rounds, following \cite{martinez2019decentralized}) is maintained exactly by the hub and broadcast in the downlink so all nodes use a consistent UCB bonus. The full pseudocode is given in Algorithm~\ref{alg:merw} (Appendix~\ref{app:pseudocode}).

\subsection{Regret Bound}

\begin{corollary}[Regret matches centralized UCB]
\label{cor:regret}
Under the \merw{} relay protocol, the expected group cumulative regret satisfies \textcolor{red}{[PLACEHOLDER: regret bound to be stated here]}, where the leading term matches a centralized UCB agent on $P$ total pulls and the dependence on wall-clock time $T$ enters through $P = \Theta(NT/(2D+1))$.
\end{corollary}

\begin{proof}[Proof sketch]
By Theorem~\ref{thm:coordinator}, the hub's information state $(\hat{n}^k, \hat{s}^k, P)$ is identical to that of a centralized UCB agent at every decision point, so the standard UCB regret analysis applies directly to the hub's pull sequence indexed by $P$. During the $2D$ communication steps of each cycle, only the hub pulls once; non-hub agents are idle and accumulate no regret. The dominant cost is therefore not an additive penalty per cycle but a reduced pull rate: each cycle of $2D+1$ wall-clock steps produces exactly one pull per agent, so over horizon $T$ the total pulls are $P = \Theta(NT/(2D+1))$. The regret bound in terms of $T$ follows by substituting this expression for $P$ into the centralized UCB guarantee.
\end{proof}

\subsection{Empirical Results}

We evaluate \merw{} on Barab\'{a}si--Albert and Erd\H{o}s--R\'{e}nyi graphs ($N=20$, $K=5$, $T=5000$, 50 runs, Figures~\ref{fig:regret_ba}--\ref{fig:regret_er}). Results on additional graph families are given in Appendix~\ref{app:regret}.

Each figure shows three panels. The first panel shows group cumulative regret $\sum_i R_i(t)$ as a function of wall-clock round $t$; communication-only rounds are included, so the curve is flat during uplink and downlink phases and steps up only when pulls occur. The second panel shows the hub's individual regret $R_{h^\star}(t)$, isolating the performance of the emergent coordinator. The third panel shows group regret as a function of total arm pulls across all agents, removing the communication overhead and measuring pure sample efficiency.

On both graph families the hub, identified from local MERW computations without any central authority, successfully coordinates the network: group regret grows logarithmically and the hub's individual regret tracks the same curve, confirming that the emergent coordinator behaves as a centralized UCB agent. The regret-vs.-pulls curve shows that sample efficiency is not degraded by the relay structure. Results on additional graph families (barbell, lollipop, chain, star, grid) are given in Appendix~\ref{app:regret}.

\begin{figure}[H]
    \centering
    \includegraphics[width=\textwidth]{results/Regret/relay_regret_ba_N20_K5_T5000.pdf}
    \caption{Cumulative regret on a Barab\'{a}si--Albert graph ($N=20$, $K=5$, $T=5000$, 50 runs). \emph{Left}: group regret vs.\ round; \emph{center}: hub regret vs.\ round; \emph{right}: group regret vs.\ total arm pulls. Shaded bands are standard error of the mean.}
    \label{fig:regret_ba}
\end{figure}

\begin{figure}[H]
    \centering
    \includegraphics[width=\textwidth]{results/Regret/relay_regret_er_N20_K5_T5000.pdf}
    \caption{Cumulative regret on an Erd\H{o}s--R\'{e}nyi graph ($N=20$, $K=5$, $T=5000$, 50 runs). Panel layout as in Figure~\ref{fig:regret_ba}.}
    \label{fig:regret_er}
\end{figure}

\section{Best Arm Identification}
\label{sec:bai}

\subsection{Setting}

In the BAI setting the horizon $T$ is not fixed. Instead, given a confidence parameter $\delta \in (0,1)$, the goal is to identify the best arm $a^\star = \arg\max_k \mu_k$ with probability at least $1-\delta$ while minimizing the total number of arm pulls across all agents.

\subsection{\merw-BAI: Hillel Elimination over the MERW Tree}

\merw-BAI adapts Hillel's successive elimination to the tree topology induced by MERW. The pull schedule and elimination rule are identical to Hillel's Algorithm 3; the difference is the communication model. Rather than all-to-all broadcast, agents communicate via the MERW spanning tree: local sums travel hop-by-hop up to the hub (uplink), the hub computes the elimination, and the surviving set $S_r$ together with the timing signal $D$ travels back down (downlink). The timing signal $D$ ensures all nodes restart the next round simultaneously, regardless of their depth in the tree.

\paragraph{Round structure.} Each elimination round $r$ has four phases:

\begin{enumerate}
    \item \textbf{Pull phase} ($L_r$ steps): every node independently computes $L_r = \lceil t_r - t_{r-1} \rceil$ from the pre-agreed schedule (all nodes know $N$, $K$, $\delta$, so no communication is needed). Each node pulls every arm in $S_{r-1}$ exactly $L_r$ times and accumulates local sums.

    \item \textbf{Uplink} ($D$ steps): local sums propagate hop-by-hop toward the hub, one hop per time step. A node at depth $d$ reaches the hub after $d$ steps.

    \item \textbf{Hub aggregation}: the hub sums all received packets to form the global aggregate $\hat{s}_k = \sum_{i=1}^N s_k^{(i)}$ for each surviving arm $k$. It then computes global averages $\tilde{p}_k = \hat{s}_k / (N \cdot L_r)$, identifies $\tilde{p}^\star = \max_k \tilde{p}_k$, and eliminates arms where $\tilde{p}_k < \tilde{p}^\star - \epsilon_r$.

    \item \textbf{Downlink} ($D$ steps): the hub broadcasts the new surviving set $S_r$ and the depth signal $D$ hop-by-hop back down the tree. Each node at depth $d$ receives the packet after $d$ steps and waits $D - d$ additional steps before starting the next pull phase, ensuring all nodes are synchronized.
\end{enumerate}

The total time cost per round is $L_r + 2D$ steps (pull + uplink + downlink), compared to $L_r + 1$ for Hillel (pull phase plus one all-to-all broadcast round). The sample complexity in terms of total arm pulls is identical: $N \cdot L_r \cdot |S_{r-1}|$ per round.

\subsection{Comparison with Hillel}

The two algorithms are identical in sample complexity by construction: both use the same $L_r$ schedule and the same elimination rule. The difference lies entirely in the communication model:

\begin{itemize}
    \item \textbf{Hillel}: assumes all-to-all broadcast in one round. The coordinator is not needed because every agent computes the elimination independently. Requires $O(N)$ messages per phase.
    \item \textbf{\merw-BAI}: uses the MERW tree. The hub acts as the aggregation point because it is the only node with the full global picture. Requires $O(N)$ messages per phase (each node sends one packet), but communication takes $O(D)$ time steps rather than $O(1)$.
\end{itemize}

The central claim is that the hub, which plays the role of aggregator in \merw-BAI, is \emph{identified from the graph structure itself} via decentralized MERW computation; no all-to-all communication is assumed. This makes \merw-BAI applicable to any undirected graph, at the cost of $O(D)$ extra time steps per round relative to Hillel. The sample complexity equivalence is formalized in Corollary~\ref{cor:bai}.

\begin{corollary}[BAI sample complexity matches Hillel]
\label{cor:bai}
Under the \merw-BAI protocol, with $\mathcal{A} = $ Hillel's successive elimination (Algorithm 3 of \cite{hillel2013distributed}), the total number of arm pulls to identify the best arm with probability at least $1-\delta$ is identical to that of Hillel's algorithm. The only overhead is $2D$ extra time steps per elimination round.
\end{corollary}

\begin{proof}
By Theorem~\ref{thm:coordinator}, the hub receives exact global sums $\hat{s}_k = \sum_i s_k^{(i)}$ after $D$ uplink steps, computes global averages $\tilde{p}_k = \hat{s}_k/(N \cdot L_r)$, and applies the elimination threshold $\epsilon_r$ identically to Hillel's centralized aggregator. The surviving set $S_r$ is therefore the same as in Hillel's algorithm. By Proposition~\ref{prop:broadcast} and the depth signal, $S_r$ reaches every agent within $D$ downlink steps, and all agents enter round $r+1$ simultaneously. Since $L_r$ and $\epsilon_r$ are pre-agreed and identical across agents, the total pull count $N \cdot \sum_r L_r \cdot |S_{r-1}|$ matches Hillel's exactly. The correctness guarantee $\Pr[\text{output} = a^\star] \geq 1-\delta$ follows from Hillel's analysis applied to the hub's elimination sequence.
\end{proof}

Empirically, we compare total arm pulls as a function of $N$ in Figure~\ref{fig:bai}, where both algorithms use the same graph and the same hub node (identified by MERW) as the aggregation point for a fair comparison.

\begin{figure}[H]
    \centering
    \includegraphics[width=0.85\textwidth]{results/BAI/merw_ucb_bai_ba_K10.pdf}
    \caption{BAI sample complexity on Barab\'{a}si--Albert graphs ($K=10$, $\delta=0.05$, 50 runs per $N$, fixed graph topology per $N$). \emph{Left}: total arm pulls vs.\ number of agents $N$. \emph{Right}: pulls per agent vs.\ $N$. Error bars are standard error of the mean. Both algorithms use the MERW-identified hub as aggregation point; \merw-BAI routes communication through the induced spanning tree while Hillel assumes direct all-to-all broadcast.}
    \label{fig:bai}
\end{figure}

\section{Discussion}
\label{sec:discussion}

\textcolor{red}{[TODO: limitations and comments. Key points to address:
\begin{itemize}
    \item The protocol assumes a connected, undirected, static graph. Dynamic or directed graphs are not handled.
    \item The hub is a single point of failure; no fault tolerance is provided.
    \item $D$ can be large on sparse or path-like graphs, degrading the effective information rate $N/(2D+1)$.
    \item The BAI version requires agents to know $N$ (for the $L_r$ schedule); the regret version does not.
    \item Power iteration convergence depends on the spectral gap $\lambda_2/\lambda_1$; on nearly-regular graphs MERW centrality may not concentrate sharply.
\end{itemize}]}

\section{Conclusion}
\label{sec:conclusion}

\textcolor{red}{[TODO: conclusion. Key points: MERW centrality as a tool for emergent coordination, the hub-and-tree backbone, matching centralized performance for both BAI and regret, the $O(D)$ communication overhead as the price of decentralization.]}

\appendix

\section{Additional MERW Visualizations}
\label{app:merw_viz}

\begin{figure}[H]
    \centering
    \includegraphics[width=\textwidth]{results/MERW_visualization/merw_tree_barbell_N20.pdf}
    \caption{Barbell graph ($N=20$). Panel layout as in Figure~\ref{fig:merw_viz_ba}.}
    \label{fig:merw_viz_barbell}
\end{figure}

\begin{figure}[H]
    \centering
    \includegraphics[width=\textwidth]{results/MERW_visualization/merw_tree_lollipop_N20.pdf}
    \caption{Lollipop graph ($N=20$). Panel layout as in Figure~\ref{fig:merw_viz_ba}.}
    \label{fig:merw_viz_lollipop}
\end{figure}

\begin{figure}[H]
    \centering
    \includegraphics[width=\textwidth]{results/MERW_visualization/merw_tree_chain_N10.pdf}
    \caption{Chain graph ($N=10$). Panel layout as in Figure~\ref{fig:merw_viz_ba}.}
    \label{fig:merw_viz_chain}
\end{figure}

\begin{figure}[H]
    \centering
    \includegraphics[width=\textwidth]{results/MERW_visualization/merw_tree_star_N20.pdf}
    \caption{Star graph ($N=20$). Panel layout as in Figure~\ref{fig:merw_viz_ba}.}
    \label{fig:merw_viz_star}
\end{figure}

\begin{figure}[H]
    \centering
    \includegraphics[width=\textwidth]{results/MERW_visualization/merw_tree_clustered_N24.pdf}
    \caption{Clustered graph ($N=24$). Panel layout as in Figure~\ref{fig:merw_viz_ba}.}
    \label{fig:merw_viz_clustered}
\end{figure}

\section{MERW: Construction and Localization}
\label{app:merw}

\subsection{Construction}

The Maximal Entropy Random Walk \cite{burda2009localization} is the unique stationary Markov chain on $G$ that assigns equal probability to all paths of equal length between any two fixed endpoints. It is the maximally unbiased random walk consistent with the graph structure.

Let $A \in \{0,1\}^{N \times N}$ be the adjacency matrix of $G$. By the Perron--Frobenius theorem, $A$ has a unique positive leading eigenvalue $\lambda_1$ with a corresponding strictly positive eigenvector $\psi \in \mathbb{R}^N_{>0}$, normalized so that $\|\psi\|_2 = 1$. The MERW transition matrix is
\[
S_{ij} = \frac{A_{ij}}{\lambda_1} \cdot \frac{\psi_j}{\psi_i},
\]
and its unique stationary distribution is $\pi_i = \psi_i^2$.

Unlike the standard random walk, whose stationary distribution is proportional to degree ($\pi_i^{\mathrm{SRW}} \propto d_i$), MERW concentrates its stationary mass on a small set of structurally central nodes. This localization phenomenon \cite{burda2009localization} makes $\psi_i$ a natural measure of centrality.

\subsection{Localization and Centrality}

The most striking property of MERW is the sharp concentration of $\pi_i = \psi_i^2$ on a small subset of nodes. On graphs with heterogeneous structure (e.g., Barab\'{a}si--Albert networks), the eigenvector $\psi$ is far from uniform: a handful of nodes carry nearly all the stationary mass, while the rest have negligible values. This mirrors the ground state of a quantum particle on the graph, where probability concentrates in the lowest-energy (most path-rich) region.

We use $\psi_i$ as the centrality score of node $i$. This is justified by two properties:
\begin{enumerate}
    \item $\psi_i$ encodes node $i$'s contribution to paths of \emph{all} lengths in the graph, not just its immediate neighborhood as degree does.
    \item $\psi$ emerges from a maximum-entropy derivation, giving it a clean information-theoretic foundation: it is the least biased measure of structural importance consistent with the graph.
\end{enumerate}

\section{Decentralized Power Iteration}
\label{app:poweriter}

In a multi-agent setting, no single agent has access to the full adjacency matrix $A$. We compute $\psi$ via the distributed power iteration of \cite{jelasity2007asynchronous}, in which each agent normalizes using only local gossip rather than a global reduce.

\paragraph{Protocol.}
Each agent $i$ maintains a local scalar $w_i^{(\tau)}$ initialized to $w_i^{(0)} = 1$. At each iteration $\tau = 1, 2, \ldots, \tau_{\mathrm{init}}$:
\begin{enumerate}
    \item Agent $i$ receives $w_j^{(\tau-1)}$ from all neighbors $j \in \mathcal{N}(i)$ and computes the unnormalized update
    \[
        \tilde{w}_i^{(\tau)} = \sum_{j \in \mathcal{N}(i)} w_j^{(\tau-1)} = (Aw^{(\tau-1)})_i.
    \]
    \item Agent $i$ computes its local log-growth rate $r_i^{(\tau)} = \log(\tilde{w}_i^{(\tau)} / w_i^{(\tau-1)})$, which approximates $\log \lambda_1$ once the iteration has converged.
    \item Agents gossip $r_i^{(\tau)}$ by pairwise averaging with random neighbors for $g$ rounds. After gossip, each $r_i$ approximates the geometric mean growth rate $\frac{1}{N}\sum_j \log(\tilde{w}_j^{(\tau)} / w_j^{(\tau-1)}) \approx \log \lambda_1$.
    \item Each agent normalizes locally: $w_i^{(\tau)} = \tilde{w}_i^{(\tau)} / \exp(r_i^{(\tau)})$.
\end{enumerate}
Iteration stops when $\max_i |w_i^{(\tau)} - w_i^{(\tau-1)}| / |w_i^{(\tau-1)}| < \epsilon$, a condition each agent checks using only its own values. After $\tau_{\mathrm{init}}$ iterations, each agent holds $\tilde{\psi}_i \approx c \cdot \psi_i$ for some common unknown scale $c > 0$. Convergence error decays geometrically as $O\!\left((\lambda_2/\lambda_1)^{\tau_{\mathrm{init}}}\right)$.

\paragraph{Why gossip normalization is decentralized.}
The standard power iteration requires a global $\ell_2$ norm at each step, which needs an all-reduce over the entire network. The gossip approach replaces this with $g$ rounds of random pairwise averaging of log-growth rates. Each agent only communicates with one random neighbor per gossip round; no agent needs to know $N$ or the global topology. With $g = O(\log N)$ gossip rounds the approximation error is negligible \cite{jelasity2007asynchronous}.

\paragraph{Why we do not need $N$.}
The key observation is that we use $\psi$ only to determine routing directions: agent $i$ routes toward neighbor $j$ if $\tilde{\psi}_j > \tilde{\psi}_i$. This requires only the \emph{ratio} $\psi_j / \psi_i$, not the absolute value of either component. The gossip normalization produces iterates proportional to $\psi$ up to a common scale, which preserves all relative orderings and is sufficient for hub identification without ever computing $N$.

\paragraph{Cost.} Each iteration requires one neighbor-exchange round plus $g$ gossip rounds, all exchanging scalar values. The eigenvector is computed once before the bandit phase and stored locally.

\section{Max-Flooding and Routing Tree}
\label{app:flooding}

\subsection{Routing Targets}

Given the approximate eigenvector $\tilde{\psi}$, each agent $i$ selects a routing target:
\[
j_i^\star = \arg\max_{j \in \mathcal{N}(i)} \tilde{\psi}_j, \quad \text{subject to } \tilde{\psi}_{j_i^\star} > \tilde{\psi}_i.
\]
If no neighbor has strictly higher $\tilde{\psi}$ than $i$, then $i$ is a \emph{local maximum}. Each agent needs only its own $\tilde{\psi}_i$ and the values $\tilde{\psi}_j$ of its immediate neighbors to determine $j_i^\star$, a purely local computation.

\subsection{Max-Flooding to a Single Tree}

The local routing rule above may produce multiple local maxima when the graph has several structurally separated dense regions. To merge all local maxima into a single spanning tree rooted at the global maximum $h^\star = \arg\max_i \tilde{\psi}_i$, we run a short \emph{max-flooding} phase after power iteration, before the bandit phase begins.

\paragraph{Max-flooding protocol.}
Each agent $i$ maintains a flooded value $m_i^{(0)} = \tilde{\psi}_i$ and a pointer $\mathrm{via}_i$ initialized to $i$ itself. The protocol runs for as many rounds as needed, with no fixed horizon:
\begin{enumerate}
    \item Each agent $i$ broadcasts $m_i^{(\tau-1)}$ to all neighbors $j \in \mathcal{N}(i)$.
    \item Agent $i$ updates: if $m_j^{(\tau-1)} > m_i^{(\tau-1)}$ for some neighbor $j$, set $m_i^{(\tau)} \leftarrow m_j^{(\tau-1)}$ and $\mathrm{via}_i \leftarrow j$; otherwise keep $m_i^{(\tau)} \leftarrow m_i^{(\tau-1)}$.
    \item If no agent updated its value in round $\tau$ (i.e., all messages were silent), the flood has converged and the protocol terminates.
\end{enumerate}
The stopping criterion in step 3 is purely local: in a synchronous model, an agent that received no update from any neighbor and sent no update itself can conclude that the flood is over. No agent needs to know $\mathrm{diam}(G)$ or any other global property of the graph. Convergence is guaranteed in at most $\mathrm{diam}(G)$ rounds, and often earlier.

After convergence, $m_i = \max_{j \in V} \tilde{\psi}_j$ for all $i$, and $\mathrm{via}_i$ points toward the neighbor through which the global maximum was first reached. Every local maximum $\ell$ checks locally whether $m_{\mathrm{via}_\ell} > \tilde{\psi}_\ell$; if so, it re-routes by setting $j_\ell^\star \leftarrow \mathrm{via}_\ell$. The unique node $h^\star$ whose own $\tilde{\psi}_{h^\star}$ equals the global maximum never sees a better value from any neighbor, so it keeps $j_{h^\star}^\star$ undefined and becomes the root. Since each re-routing step moves strictly toward the global maximum, no cycles can form.

The result is a single directed spanning tree $\mathcal{T}$ rooted at $h^\star$, with all edges oriented from children toward the root.

\begin{remark}
Max-flooding terminates in at most $\mathrm{diam}(G)$ rounds and often fewer, each round exchanging a single scalar per edge. The stopping condition requires no global coordination: each agent decides locally based on whether its own value changed. On well-connected graphs the diameter is $O(\log N)$, making this a negligible overhead relative to the bandit horizon $T$.
\end{remark}

\subsection{Routing Tree Definitions}

The routing targets, after max-flooding, define a single directed spanning tree $\mathcal{T}$ on $V$ with unique root $h^\star$. Define:
\begin{itemize}
    \item $d_i$: the depth of node $i$ in $\mathcal{T}$, i.e., the number of hops from $i$ to $h^\star$ along the routing path.
    \item $D = \max_{i \in V} d_i$: the height of $\mathcal{T}$.
    \item $\mathrm{par}(i)$: the parent of $i$ in $\mathcal{T}$ (i.e., $j_i^\star$); undefined for $h^\star$.
    \item $\mathrm{ch}(i)$: the set of children of $i$ (nodes $j$ with $j_j^\star = i$).
\end{itemize}

\section{Relay Cycle Illustration}
\label{app:relay_cycle}

Figure~\ref{fig:relay_cycle} illustrates one full relay cycle on a depth-2 tree ($D=2$, cycle length $2D+1=5$ rounds). The hub is at the top; two intermediate nodes at depth 1 each have two leaf children at depth 2.

\begin{figure}[H]
    \centering
    \includegraphics[width=\textwidth]{results/MERW_visualization/relay_cycle.pdf}
    \caption{One relay cycle on a depth-2 tree ($N=7$, $D=2$). From left to right: $t=0$ -- all nodes pull their arm locally; $t=1$ -- every node sends its packet one hop toward the hub; $t=2$ -- intermediate nodes relay the packet to the hub; $t=3$ -- hub runs UCB and sends the updated snapshot one hop down; $t=4$ -- intermediate nodes forward the snapshot to their leaves, which schedule their next pull.}
    \label{fig:relay_cycle}
\end{figure}

\section{Full Algorithm Pseudocode}
\label{app:pseudocode}

\begin{algorithm}[H]
\caption{\merw}
\label{alg:merw}
\begin{algorithmic}[1]
\Require Graph $G$, arms $K$, horizon $T$, UCB constant $c$
\State Compute $\tilde{\psi}$ via distributed power iteration (Appendix~\ref{app:poweriter})
\State Run max-flooding until no agent updates its value; each local maximum $\ell$ with $m_{\mathrm{via}_\ell} > \tilde{\psi}_\ell$ re-routes via $\mathrm{via}_\ell$
\State Single spanning tree $\mathcal{T}$ with root $h^\star$; each node knows $\mathrm{par}(i)$, $\mathrm{ch}(i)$, $d_i$, $D$
\State Every agent $i \in V$ pulls each arm $k$ once (not counted in regret): $\hat{n}_i^k \leftarrow 1$, $\hat{s}_i^k \leftarrow r_i^k$

\State \textcolor{gray}{$\triangleright$ \textbf{Cycle kick-off}}
\ForAll{$i \in V$}
    \State $a_i \leftarrow \arg\max_k \bigl[\hat{s}_i^k/\hat{n}_i^k + \sqrt{c\log P/\hat{n}_i^k}\bigr]$;\quad pull $a_i$, observe $r_i$
\EndFor
\ForAll{$i \in V \setminus \{h^\star\}$} send $(a_i,\, r_i)$ to $\mathrm{par}(i)$ \EndFor
\State $t_{\mathrm{cycle}} \leftarrow t$;\quad $t \mathrel{+}= 1$

\While{$t < T$}
    \State \textcolor{gray}{$\triangleright$ \textbf{Hub}}
    \If{uplink packet $(a, r)$ arrives}
        \State $\hat{n}^a \mathrel{+}= 1$;\quad $\hat{s}^a \mathrel{+}= r$;\quad $t_{\mathrm{last}} \leftarrow t$
    \EndIf
    \If{uplink buffer empty}
        \State $D \leftarrow t_{\mathrm{last}} - t_{\mathrm{cycle}}$
        \State $P \mathrel{+}= |\text{uplink messages received}| + 1$ \Comment{non-hub pulls + hub pull}
        \State $a_h \leftarrow \arg\max_k \bigl[\hat{s}^k/\hat{n}^k + \sqrt{c\log P/\hat{n}^k}\bigr]$;\quad pull $a_h$;\quad $\hat{n}^{a_h} \mathrel{+}= 1$, $\hat{s}^{a_h} \mathrel{+}= r$
        \State send $(\hat{n},\; \hat{s},\; D,\; P,\; \mathrm{hops}{=}0)$ to each child in $\mathrm{ch}(h^\star)$
    \EndIf
    \State \textcolor{gray}{$\triangleright$ \textbf{Non-hub nodes}}
    \ForAll{$i \in V \setminus \{h^\star\}$}
        \If{$i$ holds an uplink packet $(a, r)$}
            \State forward $(a, r)$ to $\mathrm{par}(i)$;\quad go silent
        \ElsIf{$i$ received downlink $(\hat{n}, \hat{s}, D, \mathrm{hops})$}
            \State $\hat{n}_i \leftarrow \hat{n}$;\quad $\hat{s}_i \leftarrow \hat{s}$;\quad $P_i \leftarrow P$
            \State forward $(\hat{n},\; \hat{s},\; D,\; P,\; \mathrm{hops}{+}1)$ to each child in $\mathrm{ch}(i)$
            \State schedule pull at $t + (D - \mathrm{hops} - 1)$
        \ElsIf{scheduled pull time $= t$}
            \State $a_i \leftarrow \arg\max_k \bigl[\hat{s}_i^k/\hat{n}_i^k + \sqrt{c\log P/\hat{n}_i^k}\bigr]$;\quad pull $a_i$, observe $r_i$
            \State send $(a_i, r_i)$ to $\mathrm{par}(i)$;\quad $t_{\mathrm{cycle}} \leftarrow t$
        \Else
            \State go silent
        \EndIf
    \EndFor
    \State $t \mathrel{+}= 1$
\EndWhile
\end{algorithmic}
\end{algorithm}

\bibliographystyle{plain}
\begin{thebibliography}{9}

\bibitem{burda2009localization}
Z.~Burda, J.~Duda, J.M.~Luck, and B.~Waclaw.
\newblock Localization of the maximal entropy random walk.
\newblock \emph{Physical Review Letters}, 102(16):160602, 2009.

\bibitem{landgren2021distributed}
P.~Landgren, V.~Srivastava, and N.E.~Leonard.
\newblock Distributed cooperative decision making in multi-agent multi-armed bandits.
\newblock \emph{Automatica}, 125:109445, 2021.

\bibitem{jelasity2007asynchronous}
M.~Jelasity, G.~Canright, and K.~Eng\o-Monsen.
\newblock Asynchronous distributed power iteration with gossip-based normalization.
\newblock In \emph{European Conference on Parallel Processing}, pp.~514--525, 2007.

\bibitem{martinez2019decentralized}
D.~Mart\'inez-Rubio, V.~Kanade, and P.~Rebeschini.
\newblock Decentralized cooperative stochastic bandits.
\newblock In \emph{Advances in Neural Information Processing Systems (NeurIPS)}, 2019.

\bibitem{hillel2013distributed}
E.~Hillel, Z.~Karnin, T.~Koren, R.~Lempel, and O.~Somekh.
\newblock Distributed exploration in multi-armed bandits.
\newblock In \emph{Advances in Neural Information Processing Systems (NeurIPS)}, 2013.

\end{thebibliography}

\section{Additional Regret Results}
\label{app:regret}

Figures~\ref{fig:regret_barbell}--\ref{fig:regret_grid} show cumulative regret on the remaining graph families. The three-panel layout is identical to Figure~\ref{fig:regret_ba}: group regret vs.\ round, hub regret vs.\ round, and group regret vs.\ total arm pulls.

On the barbell graph the hub concentrates in one of the two cliques; nodes in the other clique route through the bridge edge, resulting in a larger $D$ and correspondingly slower per-round progress visible in the flat segments of the wall-clock regret curve. On the lollipop, the hub sits at the clique core and the path nodes form a deep chain ($D$ up to $N/2$), making the flat segments especially pronounced. On the chain the hub is the central node and $D = N/2$, so each cycle is long; nevertheless regret vs.\ total pulls remains well-behaved. The star is the ideal case: $D=1$, every cycle lasts 3 rounds, and the hub aggregates all observations with minimal delay. On the grid MERW selects the center node and $D$ equals the grid radius, giving moderate cycle length.

Across all graph families the regret-vs.-pulls curves are consistent, confirming that the sample efficiency of \merw{} is independent of graph topology: the relay structure adds communication overhead (visible in wall-clock regret) but does not waste observations.

\begin{figure}[H]
    \centering
    \includegraphics[width=\textwidth]{results/Regret/relay_regret_barbell_N20_K5_T5000.pdf}
    \caption{Barbell graph ($N=20$, $K=5$, $T=5000$, 50 runs).}
    \label{fig:regret_barbell}
\end{figure}

\begin{figure}[H]
    \centering
    \includegraphics[width=\textwidth]{results/Regret/relay_regret_lollipop_N20_K5_T5000.pdf}
    \caption{Lollipop graph ($N=20$, $K=5$, $T=5000$, 50 runs).}
    \label{fig:regret_lollipop}
\end{figure}

\begin{figure}[H]
    \centering
    \includegraphics[width=\textwidth]{results/Regret/relay_regret_chain_N20_K5_T5000.pdf}
    \caption{Chain graph ($N=20$, $K=5$, $T=5000$, 50 runs).}
    \label{fig:regret_chain}
\end{figure}

\begin{figure}[H]
    \centering
    \includegraphics[width=\textwidth]{results/Regret/relay_regret_star_N20_K5_T5000.pdf}
    \caption{Star graph ($N=20$, $K=5$, $T=5000$, 50 runs).}
    \label{fig:regret_star}
\end{figure}

\begin{figure}[H]
    \centering
    \includegraphics[width=\textwidth]{results/Regret/relay_regret_grid_N25_K5_T5000.pdf}
    \caption{Grid graph ($N=25$, $K=5$, $T=5000$, 50 runs).}
    \label{fig:regret_grid}
\end{figure}

\end{document}
